\label{chap:Conclusion}

This thesis presented a modular, end-to-end machine learning pipeline for
earthquake phase picking, phase association, hypocenter location, and magnitude 
estimation, deployed on the CINECA Leonardo \ac{HPC} infrastructure. The 
pipeline was evaluated over a 93-day continuous waveform archive from the 
\ac{SMINO} in Northeastern Italy, spanning 266 stations across 18 networks, and 
rigorously validated against the manually reviewed OGS reference catalog of 380 
events and 7,129 picks (4,126~P and 3,003~S).

\section{Summary of Findings}

\paragraph{Phase picking.}
A comprehensive benchmarking campaign compared PhaseNet and EQTransformer
models trained on four datasets (INSTANCE, STEAD, SCEDC, and the original
training sets). PhaseNet trained on the INSTANCE dataset achieved the best
overall balance of recall and timing precision, with P-wave recall of 0.953,
S-wave recall of 0.918, and a combined recall of 0.938 at the default
detection threshold of 0.1. Its timing accuracy was the highest of all
configurations: P-phase MAE~$=0.066$\,s (mean bias~$=-0.021$\,s) and S-phase 
MAE~$= 0.116$\,s (mean bias~$=+0.054$\,s). EQTransformer (INSTANCE) attained 
marginally higher overall recall (0.944) but exhibited systematically larger 
timing residuals, P-phase MAE~$=0.075$\,s ($+14\%$) and S-phase MAE~$=0.134$\,s 
($+16\%$), with a pronounced S-phase bias of $+0.091$\,s. This indicates that 
PhaseNet provides superior temporal precision, that is critical for downstream 
association and location.

The benchmark results revealed that domain proximity is the primary determinant 
of picker performance. The INSTANCE dataset, which includes stations and events 
from Northeastern Italy (overlapping with the OGS network), produced the most 
sensitive and precise pickers. Models trained on out of domain data showed 
clear degradation: SCEDC-trained models generated $12\times$--$22\times$ more 
proposed picks (23--30~million vs. $\sim$1--2~million) due to an inability to 
discriminate regional noise characteristics, while STEAD-trained PhaseNet 
achieved only 0.750 combined recall, the lowest among all variants. The 
threshold sweep (0.1, 0.2, 0.3) demonstrated that raising the detection 
probability threshold primarily filters low-confidence picks without degrading 
timing quality: P-phase MAE remained stable at $\sim$0.065\,s across all 
thresholds, while proposed-pick volume decreased by a factor of $\sim$5$\times$ 
from 0.1 to 0.3.

\paragraph{Phase association.}
Two association algorithms, GaMMA and PyOcto, were evaluated downstream of the
selected PhaseNet (INSTANCE) picker. At the association stage, both achieved
identical event recall of 96.1\% (365 of 380 catalog events) at threshold~0.1,
but they differed substantially in their pick-level behavior. GaMMA's
probabilistic filtering retained only 70.6\% of matched picks (combined P~$+$~S
recall), discarding $\sim$25\% of true-positive arrivals during association
while reducing the proposed-pick volume by a factor of $\sim$60$\times$ (from
$\sim$2~million to $\sim$33,000). PyOcto preserved 90.7\% of matched picks, 
retaining 94\% of all matched P arrivals and 86\% of matched S arrivals, at the 
cost of somewhat higher proposed-pick counts ($\sim$51,500).

The event time residuals revealed contrasting systematic biases: GaMMA placed
origin times $+0.484$\,s later than the OGS catalog, while PyOcto placed them
$-0.421$\,s earlier, reflecting fundamental algorithmic differences in
preliminary origin-time estimation. Spatially, GaMMA achieved slightly tighter
epicentral accuracy (median 1.10\,km vs.\ PyOcto's 1.29\,km) but both exhibited 
comparable depth uncertainties at this stage.

\paragraph{Hypocenter location.}
NonLinLoc was applied with a 1D layered velocity model and revealed a striking
divergence between the two associators. GaMMA+NonLinLoc suffered a catastrophic
recall drop from 96.1\% to 73.7\% (280 of 380 events), losing 85 catalog 
events, nearly one quarter, during location. In contrast, PyOcto+NonLinLoc
maintained 96.3\% recall (366 events), actually recovering one additional event
beyond the association stage. This disparity demonstrates that PyOcto's 
association produces events with sufficient spatial constraints for robust 
hypocenter inversion, whereas GaMMA's probabilistic associations frequently 
lack the azimuthal coverage and station density required for NonLinLoc 
convergence.

After NonLinLoc inversion, both pathways achieved comparable location quality
for their respective matched-event sets. Origin time residuals were centered
near zero with small bias ($\sim -0.06$\,s for both pathways), MAE of 0.161\,s 
(GaMMA) and 0.171\,s (PyOcto), confirming that the probabilistic inversion 
effectively re-estimates origin times independent of the associator's 
preliminary values, using the same 1D velocity model as the OGS catalog. 
Epicentral distances showed median accuracy of 0.67\,km (GaMMA) and 0.72\,km 
(PyOcto), with 90\% of events located within 2.7\,km of catalog positions. 
Depth residuals exhibited a slight positive bias ($+0.50$\,km for GaMMA, 
$+0.15$\,km for PyOcto) with standard deviations of 2.75--3.22\,km.

\paragraph{Magnitude estimation.}
Both pathways showed virtually identical magnitude residual values:
mean~$\Delta M_L = -0.193$ (GaMMA) and $-0.190$ (PyOcto), with MAE
$\approx 0.22$ and median residuals of $-0.108$ and $-0.117$, respectively.
The systematic underestimation originates in station-subset differences between
the automated and manual workflows. Above $M_L \approx 2$, catalog magnitudes 
agree within $\pm 0.3$ magnitude units. Below this threshold, a characteristic 
bimodal residual pattern emerges, with one cluster near zero and another 
showing underestimation by 0.3--0.8 units. The critical practical advantage of 
the PyOcto pathway is that its magnitude estimates cover 366 events versus
GaMMA's 280, a 31\% larger catalog, providing a more complete
magnitude--frequency distribution particularly in the $M_L < 1$ range.

\paragraph{Event-level false discovery rate.}
Across all configurations and pipeline stages, event-level \ac{FDR} remained
high, ranging from 0.656 to 0.781. At the association stage with threshold~0.1,
GaMMA and PyOcto produced \ac{FDR} values of 0.744 and 0.781, respectively,
corresponding to roughly 3--4 proposed events for every catalog-confirmed
event. These values persisted through NonLinLoc location essentially
unchanged (0.743 for GaMMA+NonLinLoc; 0.781 for PyOcto+NonLinLoc at
threshold~0.1), indicating that proposed events possess internally consistent
pick sets and plausible hypocentral solutions that survive probabilistic
inversion. Raising the detection threshold to 0.3 reduced \ac{FDR} to 0.658
(GaMMA+NonLinLoc) and 0.675 (PyOcto+NonLinLoc), demonstrating that threshold
tuning is the primary mechanism for controlling proposed-event volume.
Raising the detection probability threshold from 0.1 to 0.3 reduced the 
\ac{FDR} appreciably: from 0.743 to 0.658 for GaMMA+NonLinLoc and from 0.781 
to 0.675 for PyOcto+NonLinLoc. This improvement occurs because a higher 
threshold discards low-confidence picks, which disproportionately contribute to 
spurious event proposals. However, this \ac{FDR} reduction comes at the cost 
of recall: event-level recall decreased from 96.1\% at threshold~0.1 to 92.1\% 
at threshold~0.3, meaning that fewer catalog events are recovered as the 
threshold becomes more stringent. This trade-off between recall and \ac{FDR} 
is a fundamental characteristic of threshold-based detection systems: lowering 
the threshold maximizes sensitivity (high recall) at the expense of more false 
proposals (high \ac{FDR}), while raising it improves purity (lower \ac{FDR}) 
but risks missing genuine events (lower recall).

It is important to note that the OGS reference catalog is itself incomplete, 
particularly below $M_L \approx 1$, where many micro-earthquakes go 
unreviewed. Consequently, a substantial fraction of the proposed events may 
represent genuine uncataloged micro-seismicity rather than false detections. 
This interpretation is supported by the spatial distribution of proposed 
events, which are concentrated in patterns consistent with known seismogenic 
zones. Nonetheless, the raw automated catalog requires post-processing 
quality-control filters, such as minimum station count, maximum azimuthal 
gap, and travel-time residual threshold, before operational use to reduce 
the effective \ac{FDR} to acceptable levels.

\paragraph{Missing events.}
The analysis of missed events revealed fundamentally different failure modes
for the two associators. GaMMA+NonLinLoc missed 100 events. PyOcto+NonLinLoc 
missed only 14 events. Spatially, PyOcto's missed events clustered at the
periphery of the bulletin area (red polygon), where station density is 
insufficient for robust location, while GaMMA's missed events were scattered 
throughout the monitoring area.

\section{Implications for Operational Seismology}
The results demonstrate that publicly available deep learning models, accessed
through the SeisBench framework (\cite{SeisBench}), can achieve recall levels
approaching, and in the best configuration (PyOcto+NonLinLoc, 96.3\%)
effectively matching, those of expert human analysts when applied to a regional
network with an appropriately matched training dataset. The 93-day evaluation
period encompassed the full range of seismicity recorded by the OGS network,
from micro-earthquakes ($M_L < 0$) to moderate events ($M_L > 3$), providing
a statistically robust assessment of pipeline reliability.

The choice of training dataset emerged as the single most important
determinant of system performance, outweighing the choice of model
architecture (PhaseNet vs.\ EQTransformer) or detection threshold. The
INSTANCE dataset's inclusion of Northeastern Italian stations provided the
domain alignment necessary for high recall and precise timing, while
out-of-domain datasets (SCEDC, STEAD) demonstrated that geographic and
tectonic mismatch can degrade performance by 10--20 percentage points in
recall or inflate proposed-pick volumes by more than an order of magnitude.
This finding underscores the practical importance of regional training data
for operational deployments.

The comparison of GaMMA and PyOcto revealed that the choice of associator
has a far greater impact on end-to-end pipeline performance than the choice
of picker or detection threshold. Despite identical association-stage recall,
PyOcto preserved 30\% more catalog events through the location stage (366 vs.\ 
280), demonstrating that spatial consistency in association is essential for 
robust hypocenter determination.

The high event-level \ac{FDR} across all configurations (0.656--0.781)
indicates that the raw automated output contains a substantial number of
proposed events that do not match the OGS catalog. At threshold~0.1, the
best-performing pathway (PyOcto+NonLinLoc) produces $\sim$3.6 proposed events
for every MH detection (\ac{FDR}~0.781). However, the persistence of these
proposed events through NonLinLoc inversion, combined with their spatial
clustering in patterns consistent with known seismicity, suggests that many
represent genuine micro-seismicity. Nonetheless, without systematic validation
of proposed events, the raw automated output would require post-processing
filters to achieve the reliability expected in operational seismic monitoring.

The \ac{HPC} deployment on the Leonardo supercomputer at CINECA enables both
near-real-time analysis and large-scale historical reprocessing on the same
platform. The hybrid parallelization strategy (MPI for task distribution,
CUDA for GPU-accelerated inference, OpenMP for thread-level parallelism)
provides computational flexibility that is essential for processing
terabyte-scale continuous waveform archives. The Makefile-based workflow
automation ensures reproducibility and allows individual pipeline components
(picker, associator, locator) to be upgraded or replaced independently as
improved models become available, without disrupting the overall workflow.

\subsection{Operational Recommendations}
Based on the findings of this study, the following practical guidelines are
proposed for deploying the pipeline in operational settings:
\begin{itemize}
  \item \textbf{Picker}: Use PhaseNet trained on INSTANCE as the default
  picker; validate new deployments with time-residual diagnostics and
  recall/\ac{FDR} distributions.
  \item \textbf{Associator}: Prefer PyOcto for the area under investigation.
  GaMMA should be tuned to improve its performance for the region.
  \item \textbf{Thresholds}: Start at 0.2--0.3 for P and S detection
  probabilities; adjust per-station or per-region if systematic
  recall/\ac{FDR} patterns appear in quality-control dashboards.
  \item \textbf{Quality-control gates}: Enforce minimum station count, maximum
  per-station residual, and P/S ordering checks before publishing events.
  \item \textbf{Monitoring}: Track key performance indicators (recall,
  \ac{FDR}, processing lead time) and re-tune monthly or after major
  configuration changes.
\end{itemize}

\section{Limitations and Future Work}
Several limitations of the current study and directions for future development
should be noted:
\begin{itemize}
  \item \textbf{3D velocity models.} All locations in this study were computed
  with a 1D layered velocity model. Replacing the 1D model with a regional 3D velocity model is
  expected to improve event characterization,
  particularly for shallow events and those near major structural boundaries.

  \item \textbf{Threshold optimization.} This study evaluated uniform detection
  thresholds (0.1, 0.2, 0.3) applied identically across all stations and
  phases. Systematic exploration of per-station or per-region adaptive
  thresholds, informed by station noise characteristics and local seismicity
  rates, could improve the recall/\ac{FDR} balance beyond the uniform
  settings tested here.

  \item \textbf{Magnitude calibration.} The bimodal residual pattern observed
  below $M_L \approx 2$, with systematic underestimation of 0.3--0.8 magnitude
  units for a subset of events, requires further investigation.

  \item \textbf{Additional associators.} Evaluating alternative association
  algorithms (e.g., REAL (\cite{REAL})) is recommended.

  \item \textbf{Proposed-event validation.} The 810--1,304 proposed events
  absent from the OGS catalog require systematic validation. Many likely 
  represent genuine uncataloged micro-seismicity, and quality-control criteria 
  such as minimum station count, maximum azimuthal gap, and travel-time
  residual thresholds should be applied.
\end{itemize}

\section{Concluding Remarks}
This thesis demonstrated that a modular ML pipeline, combining PhaseNet
(INSTANCE) for phase picking, PyOcto for phase association, NonLinLoc for
hypocenter location, and local magnitude estimation, can recover 96.3\% of a
manually reviewed regional earthquake catalog with sub-kilometer epicentral
accuracy (median 0.72\,km), sub-second origin time precision (MAE 0.171\,s),
and magnitude agreement within $\pm 0.22$ magnitude units. These results
establish a quantitative baseline for automated seismic monitoring in
Northeastern Italy and demonstrate that deep-learning-based pipelines deployed
on HPC infrastructure are ready for operational integration into regional
seismic surveillance systems.